\documentclass[12pt,a4paper,openany,oneside]{book}

% --- Paquetes ---
\usepackage[utf8]{inputenc}
\usepackage[spanish, es-tabla]{babel}
\usepackage{amsmath, amsfonts, amssymb}
\usepackage{graphicx}
\usepackage{geometry}
\usepackage{hyperref}
\usepackage{booktabs}
\usepackage{fancyhdr}
\usepackage{listings}
\usepackage{xcolor}
\usepackage{tikz}
\usepackage{enumitem}
\usepackage{array}
\usetikzlibrary{arrows.meta, positioning, shapes.geometric}

\geometry{margin=2.5cm}

% --- Estilo de Código Python ---
\definecolor{codegreen}{rgb}{0,0.6,0}
\definecolor{codegray}{rgb}{0.5,0.5,0.5}
\definecolor{codepurple}{rgb}{0.58,0,0.82}
\definecolor{backcolour}{rgb}{0.95,0.95,0.92}

\lstset{
    language=Python,
    backgroundcolor=\color{backcolour},   
    commentstyle=\color{codegreen},
    keywordstyle=\color{blue},
    numberstyle=\tiny\color{codegray},
    stringstyle=\color{codepurple},
    basicstyle=\ttfamily\small,
    breaklines=true,
    frame=single,
    showstringspaces=false
}

% --- Cabecera y Pie de Página ---
\pagestyle{fancy}
\fancyhf{}
\renewcommand{\headrulewidth}{0.4pt}
\renewcommand{\footrulewidth}{0.4pt}
\fancyhead[L]{\small Carlos César Sánchez Coronel}
\fancyhead[R]{\small Procesamiento de Lenguaje Natural (NLP)}
\fancyfoot[L]{\small Carlos César Sánchez Coronel}
\fancyfoot[C]{\thepage}

% --- Portada ---
\title{
    \textbf{Sílabo del Curso} \\ 
    \Huge \textbf{PROCESAMIENTO DE LENGUAJE NATURAL (NLP)} \\
    \Large \textsc{para Inteligencia Artificial} \\[1cm]
    \normalsize De los fundamentos lingüísticos a los modelos generativos
}
\author{
    \small Por \\ \vspace{0.2cm}
    \Large \textsc{Carlos César Sánchez Coronel}
}
\date{2026}

\begin{document}

\maketitle
\tableofcontents

% ------------------------------------------------------------------
% PRESENTACIÓN DEL CURSO
% ------------------------------------------------------------------
\chapter*{Presentación del Curso}
\addcontentsline{toc}{chapter}{Presentación del Curso}

Este curso ha sido diseñado para futuros ingenieros en inteligencia artificial que deseen dominar el Procesamiento de Lenguaje Natural (NLP) desde una perspectiva práctica y orientada a la industria. A lo largo de 14 semanas, exploraremos desde las técnicas clásicas de preprocesamiento hasta la arquitectura de transformers y los modelos generativos, con un enfoque en las preguntas y habilidades que se esperan en entrevistas técnicas.

Cada tema será abordado con teoría sólida, implementaciones en Python usando librerías estándar (NLTK, spaCy, Hugging Face, PyTorch) y ejemplos de casos reales como chatbots, análisis de sentimiento, sistemas de búsqueda semántica y asistentes virtuales. Al finalizar, el estudiante estará preparado para enfrentar desafíos de modelado de lenguaje y comprender el ciclo completo de un proyecto de NLP, desde la recolección de datos hasta el despliegue en producción.

\vspace{0.5cm}
\textbf{Estructura del curso (14 semanas):}
\begin{itemize}
    \item[$\square$] Semanas 1-3: Fundamentos y preprocesamiento
    \item[$\square$] Semanas 4-6: Representación del lenguaje y modelos clásicos
    \item[$\square$] Semanas 7-9: Modelos de deep learning para secuencias
    \item[$\square$] Semanas 10-12: Transformers y modelos pre-entrenados
    \item[$\square$] Semanas 13-14: Aplicaciones avanzadas y despliegue
\end{itemize}

% ------------------------------------------------------------------
% SEMANA 1: INTRODUCCIÓN A NLP Y FUNDAMENTOS LINGÜÍSTICOS
% ------------------------------------------------------------------
\chapter{Semana 1: Introducción a NLP y Fundamentos Lingüísticos}

\section{Logro de la sesión}
Comprender el alcance del NLP, sus aplicaciones en la industria y los conceptos lingüísticos básicos que subyacen al procesamiento del lenguaje.

\subsection{Conceptos}
\begin{itemize}
    \item Definición de NLP: objetivos, desafíos y áreas de aplicación.
    \item Niveles del lenguaje: morfología, sintaxis, semántica, pragmática.
    \item Part-of-Speech (POS) tagging y parsing.
    \item Corpus y datasets anotados.
    \item Aplicaciones reales: chatbots, análisis de sentimiento, traducción automática, búsqueda semántica.
\end{itemize}

\subsection{Aplicaciones en el mundo laboral y entrevistas}
\begin{itemize}
    \item Pregunta típica: "¿Qué es NLP y dónde se usa?" Se espera mencionar casos concretos como atención al cliente, detección de fraudes en textos financieros, análisis de redes sociales.
    \item Entender la diferencia entre sintaxis (estructura) y semántica (significado).
\end{itemize}

\subsection{Python}
\begin{itemize}
    \item Introducción a NLTK y spaCy: carga de pipelines, exploración de corpora.
    \item Visualización de árboles de dependencia con spaCy.
\end{itemize}

% ------------------------------------------------------------------
% SEMANA 2: PREPROCESAMIENTO DE TEXTO
% ------------------------------------------------------------------
\chapter{Semana 2: Preprocesamiento de Texto}

\section{Logro de la sesión}
Aplicar técnicas de limpieza y normalización de texto para preparar datos para modelos de NLP.

\subsection{Conceptos}
\begin{itemize}
    \item Tokenización: palabras, oraciones, subpalabras (BPE, WordPiece, Unigram).
    \item Normalización: minúsculas, eliminación de caracteres especiales, manejo de URLs y emojis.
    \item Eliminación de stopwords.
    \item Stemming vs. Lemmatization: diferencias y cuándo usar cada una.
    \item Manejo de expresiones regulares (regex) para limpieza avanzada.
    \item Desafíos en NLP: Out-of-Vocabulary (OOV) words y cómo los métodos de subpalabras los resuelven.
\end{itemize}

\subsection{Aplicaciones y entrevistas}
\begin{itemize}
    \item Pregunta clave: "¿Cuándo usar stemming y cuándo lematización?"
    \item Conocimiento de algoritmos de tokenización como BPE (usado en GPT) y WordPiece (usado en BERT).
\end{itemize}

\subsection{Python}
\begin{itemize}
    \item Uso de NLTK y spaCy para tokenización, lematización y eliminación de stopwords.
    \item Práctica con regex para limpieza de texto ruidoso.
    \item Implementación simple de BPE con librerías como \texttt{tokenizers}.
\end{itemize}

% ------------------------------------------------------------------
% SEMANA 3: REPRESENTACIÓN DEL TEXTO: BOW, TF-IDF Y N-GRAMAS
% ------------------------------------------------------------------
\chapter{Semana 3: Representación del Texto: BOW, TF-IDF y N-gramas}

\section{Logro de la sesión}
Transformar texto en vectores numéricos utilizando técnicas clásicas de representación.

\subsection{Conceptos}
\begin{itemize}
    \item Bag of Words (BOW): frecuencia de términos, limitaciones (pérdida de orden y semántica).
    \item N-gramas: capturar contexto local.
    \item TF-IDF (Term Frequency-Inverse Document Frequency): ponderación de términos según su importancia.
    \item Limitaciones de las representaciones dispersas (sparse representations).
\end{itemize}

\subsection{Aplicaciones y entrevistas}
\begin{itemize}
    \item Pregunta frecuente: "Compara Bag of Words, TF-IDF y embeddings. ¿Cuándo usarías cada uno?"
    \item Saber que TF-IDF penaliza palabras muy comunes (como artículos) y resalta términos distintivos.
\end{itemize}

\subsection{Python}
\begin{itemize}
    \item Implementación con \texttt{sklearn.feature\_extraction.text.CountVectorizer} y \texttt{TfidfVectorizer}.
    \item Análisis de documentos con estas representaciones.
\end{itemize}

% ------------------------------------------------------------------
% SEMANA 4: WORD EMBEDDINGS: WORD2VEC Y GLOVE
% ------------------------------------------------------------------
\chapter{Semana 4: Word Embeddings: Word2Vec y GloVe}

\section{Logro de la sesión}
Comprender el concepto de embeddings densos y entrenar/visualizar word embeddings.

\subsection{Conceptos}
\begin{itemize}
    \item Limitaciones de las representaciones dispersas: maldición de la dimensionalidad, falta de semántica.
    \item Word2Vec: modelos Skip-gram y CBOW, entrenamiento con ventanas de contexto.
    \item GloVe (Global Vectors): factorización de matrices de co-ocurrencia.
    \item Propiedades de los embeddings: relaciones semánticas (rey - hombre + mujer ≈ reina).
    \item Embeddings contextuales vs. estáticos.
\end{itemize}

\subsection{Aplicaciones y entrevistas}
\begin{itemize}
    \item Pregunta técnica: "Explica la diferencia entre Skip-gram y CBOW".
    \item Saber interpretar la analogía de embeddings en entrevistas.
\end{itemize}

\subsection{Python}
\begin{itemize}
    \item Entrenar Word2Vec con \texttt{gensim} sobre un corpus pequeño.
    \item Visualizar embeddings con PCA/t-SNE y matplotlib.
    \item Cargar embeddings pre-entrenados de GloVe y explorar similitudes.
\end{itemize}

% ------------------------------------------------------------------
% SEMANA 5: MODELOS CLÁSICOS PARA NLP (NAIVE BAYES, SVM)
% ------------------------------------------------------------------
\chapter{Semana 5: Modelos Clásicos para NLP (Naive Bayes, SVM)}

\section{Logro de la sesión}
Aplicar algoritmos de machine learning clásicos a tareas de NLP como clasificación de texto.

\subsection{Conceptos}
\begin{itemize}
    \item Naive Bayes para clasificación de texto: supuesto de independencia, cálculo de probabilidades.
    \item Support Vector Machines (SVM) con kernels lineales para textos.
    \item Ventajas y limitaciones de modelos clásicos frente a deep learning.
\end{itemize}

\subsection{Aplicaciones y entrevistas}
\begin{itemize}
    \item Pregunta típica: "Menciona dos algoritmos clásicos para clasificación de textos".
    \item Saber que Naive Bayes es rápido y funciona bien con datos pequeños, pero asume independencia de características.
\end{itemize}

\subsection{Python}
\begin{itemize}
    \item Implementación de clasificador Naive Bayes con \texttt{sklearn} sobre dataset de noticias (20 Newsgroups).
    \item Comparación con SVM lineal.
    \item Evaluación con métricas de clasificación (accuracy, F1-score).
\end{itemize}

% ------------------------------------------------------------------
% SEMANA 6: REDES NEURONALES PARA NLP Y EMBEDDINGS ENTRENABLES
% ------------------------------------------------------------------
\chapter{Semana 6: Redes Neuronales para NLP y Embeddings Entrenables}

\section{Logro de la sesión}
Construir redes neuronales simples para tareas de NLP usando PyTorch.

\subsection{Conceptos}
\begin{itemize}
    \item De los embeddings estáticos a los embeddings entrenables.
    \item Capa de embedding en PyTorch.
    \item Redes feed-forward para clasificación de texto (promedio de embeddings + MLP).
    \item Función de pérdida: Cross-Entropy.
    \item Regularización en NLP: dropout, weight decay.
\end{itemize}

\subsection{Aplicaciones y entrevistas}
\begin{itemize}
    \item Pregunta: "¿Cómo representarías una oración para ingresarla a una red neuronal?" (respuesta: promedio de embeddings o concatenación).
    \item Saber que los embeddings entrenables se ajustan a la tarea específica.
\end{itemize}

\subsection{Python}
\begin{itemize}
    \item Implementación de un clasificador simple con embeddings entrenables en PyTorch.
    \item Entrenamiento sobre dataset de IMDb para análisis de sentimiento.
    \item Comparación con embeddings pre-entrenados estáticos.
\end{itemize}

% ------------------------------------------------------------------
% SEMANA 7: REDES RECURRENTES (RNN, LSTM, GRU) PARA NLP
% ------------------------------------------------------------------
\chapter{Semana 7: Redes Recurrentes (RNN, LSTM, GRU) para NLP}

\section{Logro de la sesión}
Modelar secuencias de texto utilizando redes recurrentes y entender sus limitaciones.

\subsection{Conceptos}
\begin{itemize}
    \item RNN (Recurrent Neural Networks): cómo procesan secuencias, estado oculto.
    \item Problema del gradiente desvaneciente (vanishing gradient).
    \item LSTM (Long Short-Term Memory): celdas de memoria, puertas (input, forget, output).
    \item GRU (Gated Recurrent Unit): versión simplificada de LSTM.
\end{itemize}

\subsection{Aplicaciones y entrevistas}
\begin{itemize}
    \item Pregunta avanzada: "Explica por qué LSTM funciona mejor que RNN vanilla para secuencias largas".
    \item Saber que LSTM fue el estado del arte antes de los transformers.
\end{itemize}

\subsection{Python}
\begin{itemize}
    \item Implementación de LSTM para clasificación de sentimiento con PyTorch.
    \item Comparación de rendimiento entre RNN simple y LSTM.
    \item Uso de embeddings pre-entrenados con LSTM.
\end{itemize}

% ------------------------------------------------------------------
% SEMANA 8: MODELOS SEQ2SEQ Y ATENCIÓN
% ------------------------------------------------------------------
\chapter{Semana 8: Modelos Seq2Seq y Atención}

\section{Logro de la sesión}
Construir modelos de secuencia a secuencia (Seq2Seq) para tareas como traducción automática y summarization.

\subsection{Conceptos}
\begin{itemize}
    \item Arquitectura Encoder-Decoder.
    \item Context vector fijo y su limitación.
    \item Mecanismo de atención: alineación entre palabras de entrada y salida.
    \item Tipos de atención: global vs. local.
    \item Métricas para traducción: BLEU (Bilingual Evaluation Understudy).
\end{itemize}

\subsection{Aplicaciones y entrevistas}
\begin{itemize}
    \item Pregunta fundamental: "¿Cómo funciona el mecanismo de atención?"
    \item La atención es la base de los transformers; entenderla es clave.
    \item Saber que BLEU mide la similitud entre traducción automática y referencia.
\end{itemize}

\subsection{Python}
\begin{itemize}
    \item Implementación de un modelo Seq2Seq con atención usando PyTorch.
    \item Entrenamiento para tarea de traducción simple (ej. inglés-español con dataset pequeño).
    \item Evaluación con BLEU usando \texttt{sacrebleu}.
\end{itemize}

% ------------------------------------------------------------------
% SEMANA 9: TRANSFORMERS - ARQUITECTURA Y ATENCIÓN MULTI-HEAD
% ------------------------------------------------------------------
\chapter{Semana 9: Transformers - Arquitectura y Atención Multi-Head}

\section{Logro de la sesión}
Comprender la arquitectura del transformer y su revolución en NLP.

\subsection{Conceptos}
\begin{itemize}
    \item El paper "Attention is All You Need".
    \item Arquitectura del transformer: encoder, decoder, multi-head attention, positional encoding, feed-forward layers.
    \item Self-attention y atención cruzada (cross-attention).
    \item Ventajas sobre RNN: paralelización, captura de dependencias largas.
\end{itemize}

\subsection{Aplicaciones y entrevistas}
\begin{itemize}
    \item Pregunta estrella: "Explica cómo funciona el transformer y qué problema resuelve".
    \item Saber qué es positional encoding y por qué es necesario.
\end{itemize}

\subsection{Python}
\begin{itemize}
    \item Implementación de una capa de self-attention desde cero con PyTorch.
    \item Exploración de la librería Hugging Face Transformers.
    \item Uso de modelos pre-entrenados para inferencia.
\end{itemize}

% ------------------------------------------------------------------
% SEMANA 10: MODELOS PRE-ENTRENADOS (BERT, GPT) Y FINE-TUNING
% ------------------------------------------------------------------
\chapter{Semana 10: Modelos Pre-entrenados (BERT, GPT) y Fine-tuning}

\section{Logro de la sesión}
Utilizar y fine-tunar modelos pre-entrenados como BERT y GPT para tareas específicas.

\subsection{Conceptos}
\begin{itemize}
    \item BERT (Bidirectional Encoder Representations from Transformers): entrenamiento con MLM (Masked Language Model) y NSP (Next Sentence Prediction).
    \item GPT (Generative Pre-trained Transformer): entrenamiento autorregresivo.
    \item Fine-tuning vs. feature-based approaches.
    \item Tokenizadores de subpalabras (WordPiece, BPE).
    \item Métricas para clasificación: accuracy, F1, AUC.
\end{itemize}

\subsection{Aplicaciones y entrevistas}
\begin{itemize}
    \item Pregunta común: "¿Cuál es la diferencia entre BERT y GPT?"
    \item Saber cuándo usar fine-tuning completo y cuándo usar adaptadores (LoRA).
\end{itemize}

\subsection{Python}
\begin{itemize}
    \item Fine-tuning de BERT para clasificación de textos usando Hugging Face.
    \item Uso de GPT para generación de texto.
    \item Implementación de early stopping y guardado de checkpoints.
\end{itemize}

% ------------------------------------------------------------------
% SEMANA 11: TAREAS AVANZADAS: NER, QUESTION ANSWERING, SUMMARIZATION
% ------------------------------------------------------------------
\chapter{Semana 11: Tareas Avanzadas: NER, Question Answering, Summarization}

\section{Logro de la sesión}
Aplicar modelos pre-entrenados a tareas específicas de NLP como Named Entity Recognition (NER) y extracción de respuestas.

\subsection{Conceptos}
\begin{itemize}
    \item Named Entity Recognition (NER): etiquetado de entidades (personas, organizaciones, lugares).
    \item Question Answering (QA): modelos extractivos (SQuAD) y generativos.
    \item Summarization: extractivo vs. abstractivo.
    \item Métricas de evaluación: F1 para NER, exact match para QA, ROUGE para summarization.
\end{itemize}

\subsection{Aplicaciones y entrevistas}
\begin{itemize}
    \item Pregunta: "¿Qué métrica usarías para evaluar un sistema de NER?" (F1-score por entidad).
    \item Saber que ROUGE mide solapamiento de n-gramas entre resumen generado y referencia.
\end{itemize}

\subsection{Python}
\begin{itemize}
    \item Fine-tuning de BERT para NER con Hugging Face.
    \item Uso de pipelines de Hugging Face para QA y summarization.
    \item Evaluación con métricas específicas.
\end{itemize}

% ------------------------------------------------------------------
% SEMANA 12: RAG (RETRIEVAL-AUGMENTED GENERATION) Y BASES DE DATOS VECTORIALES
% ------------------------------------------------------------------
\chapter{Semana 12: RAG (Retrieval-Augmented Generation) y Bases de Datos Vectoriales}

\section{Logro de la sesión}
Diseñar sistemas que combinan recuperación de información con generación de lenguaje (RAG).

\subsection{Conceptos}
\begin{itemize}
    \item Limitaciones de los LLMs: conocimiento desactualizado, alucinaciones.
    \item RAG (Retrieval-Augmented Generation): recuperar información relevante de una base de conocimiento y pasarla al LLM como contexto.
    \item Embeddings para búsqueda semántica.
    \item Bases de datos vectoriales: FAISS, Pinecone, Weaviate.
    \item Evaluación de sistemas RAG: faithfulness, relevance.
\end{itemize}

\subsection{Aplicaciones y entrevistas}
\begin{itemize}
    \item Pregunta clave: "¿Qué es RAG y cuándo lo usarías?"
    \item Saber que RAG reduce alucinaciones y permite usar conocimiento actualizado sin reentrenar.
\end{itemize}

\subsection{Python}
\begin{itemize}
    \item Construcción de un pipeline RAG: cargar documentos, crear embeddings con Sentence Transformers, almacenar en FAISS, recuperar y generar con GPT.
    \item Implementación de un chatbot que responda sobre documentos propios.
\end{itemize}

% ------------------------------------------------------------------
% SEMANA 13: PROMPT ENGINEERING Y OPTIMIZACIÓN DE MODELOS
% ------------------------------------------------------------------
\chapter{Semana 13: Prompt Engineering y Optimización de Modelos}

\section{Logro de la sesión}
Dominar técnicas de prompt engineering y métodos de optimización como LoRA y cuantización.

\subsection{Conceptos}
\begin{itemize}
    \item Prompt engineering: zero-shot, few-shot, chain-of-thought.
    \item In-context learning.
    \item Fine-tuning eficiente: LoRA (Low-Rank Adaptation), adaptadores.
    \item Cuantización: reducción de precisión para inferencia más rápida (INT8, FP16).
    \item Exportación de modelos a ONNX.
\end{itemize}

\subsection{Aplicaciones y entrevistas}
\begin{itemize}
    \item Pregunta avanzada: "¿Qué es LoRA y por qué es útil?"
    \item Saber que la cuantización reduce el tamaño del modelo y acelera inferencia a costa de pequeña pérdida de precisión.
\end{itemize}

\subsection{Python}
\begin{itemize}
    \item Experimentación con prompts para GPT usando librerías como \texttt{openai} o modelos locales.
    \item Fine-tuning con LoRA usando \texttt{peft} y \texttt{transformers}.
    \item Cuantización con \texttt{torch.quantization} u ONNX.
\end{itemize}

% ------------------------------------------------------------------
% SEMANA 14: DESPLIEGUE DE MODELOS NLP (MLOPS)
% ------------------------------------------------------------------
\chapter{Semana 14: Despliegue de Modelos NLP (MLOps)}

\section{Logro de la sesión}
Llevar un modelo NLP a producción, crear APIs y monitorear su rendimiento.

\subsection{Conceptos}
\begin{itemize}
    \item Separación entre entrenamiento e inferencia.
    \item Serialización de modelos: pickle, joblib, ONNX, TorchScript.
    \item Creación de APIs REST con Flask o FastAPI.
    \item Contenerización con Docker.
    \item Monitoreo de modelos en producción: data drift, concept drift, latencia, throughput.
    \item Versionamiento de modelos con DVC o Hugging Face Hub.
\end{itemize}

\subsection{Aplicaciones y entrevistas}
\begin{itemize}
    \item Pregunta clave: "¿Cómo desplegarías un modelo NLP en producción?"
    \item Saber que se debe monitorear la deriva de datos y la latencia.
\end{itemize}

\subsection{Python}
\begin{itemize}
    \item Creación de una API con FastAPI que carga un modelo fine-tuneado y recibe textos para clasificar.
    \item Dockerización de la aplicación.
    \item Pruebas de carga y monitoreo básico.
\end{itemize}

% ------------------------------------------------------------------
% METODOLOGÍA Y EVALUACIÓN
% ------------------------------------------------------------------
\chapter{Metodología y Evaluación}

\section{Estrategias metodológicas}
\begin{itemize}
    \item \textbf{Clases síncronas:} Exposición conceptual con diapositivas, programación en vivo con Python utilizando cuadernos de Jupyter y Hugging Face.
    \item \textbf{Trabajo asíncrono:} Lectura de papers seminales ("Attention is All You Need"), ejercicios prácticos y proyectos semanales.
    \item \textbf{Aprendizaje activo:} Simulacros de entrevistas técnicas, discusión de casos reales, retos de código.
    \item \textbf{Recursos:} Plataforma virtual, Zoom, grabaciones, foros y repositorio GitHub.
\end{itemize}

\section{Materiales educativos}
\begin{itemize}
    \item Cuadernos de Jupyter con ejemplos y ejercicios.
    \item Acceso a modelos en Hugging Face y datasets.
    \item Bibliografía: Jurafsky \& Martin (Speech and Language Processing), papers originales.
    \item Google Colab Pro para GPU.
\end{itemize}

\section{Evaluación}
\begin{table}[h]
\centering
\begin{tabular}{lc}
\toprule
\textbf{Ítem} & \textbf{Ponderación} \\
\midrule
Evaluaciones continuas (EC) & 40\% \\
Examen parcial (EP) - Semana 7 & 30\% \\
Examen final (EF) - Semana 14 & 30\% \\
\midrule
\textbf{Promedio final (PF)} & PF = 0.4 \times EC + 0.3 \times EP + 0.3 \times EF \\
\bottomrule
\end{tabular}
\caption{Sistema de evaluación}
\end{table}
\begin{itemize}
    \item \textbf{Evaluaciones continuas:} Controles de lectura, entrega de ejercicios de programación, participación en foros y quiz semanales.
    \item \textbf{Examen parcial:} Cubre semanas 1-7 (fundamentos, word embeddings, RNN, atención).
    \item \textbf{Examen final:} Proyecto integrador (ej. construir un chatbot con RAG y desplegarlo como API).
    \item \textbf{Examen sustitutorio:} Semana 15, reemplaza la nota más baja entre parcial y final.
\end{itemize}

% ------------------------------------------------------------------
% BIBLIOGRAFÍA
% ------------------------------------------------------------------
\chapter*{Bibliografía}
\addcontentsline{toc}{chapter}{Bibliografía}

\begin{itemize}
    \item Jurafsky, D., \& Martin, J. H. (2023). \textit{Speech and Language Processing} (3rd ed. draft). Stanford University.
    \item Vaswani, A., et al. (2017). \textit{Attention is All You Need}. NeurIPS.
    \item Devlin, J., et al. (2019). \textit{BERT: Pre-training of Deep Bidirectional Transformers for Language Understanding}. NAACL.
    \item Radford, A., et al. (2018). \textit{Improving Language Understanding by Generative Pre-Training} (GPT). OpenAI.
    \item Lewis, P., et al. (2020). \textit{Retrieval-Augmented Generation for Knowledge-Intensive NLP Tasks}. NeurIPS.
    \item Documentación oficial de Hugging Face Transformers, spaCy, NLTK, PyTorch.
\end{itemize}

\vspace{1cm}
\begin{flushright}
\textit{Elaborado por:} Mg. Carlos César Sánchez Coronel \\
\textit{Aprobado por:} Dirección de Ingeniería de Inteligencia Artificial \\
\textit{Fecha:} 29/03/2024
\end{flushright}

\end{document}